\documentclass[11pt]{article}
\usepackage[margin=1in]{geometry}
\usepackage{array}
\usepackage{booktabs}
\usepackage{longtable}
\usepackage[T1]{fontenc}
\usepackage[utf8]{inputenc}
\usepackage{hyperref}
\usepackage{xcolor}

\title{ Point-to-Point Response Letter }
\author{ review-master Export }
\date{}

\begin{document}
\maketitle

\section*{ Point-to-Point Response Letter }
% This LaTeX preview is rendered from review-master.db.

\subsection*{ reviewer\_1 }
\begin{longtable}{|p{0.22\textwidth}|p{0.20\textwidth}|p{0.24\textwidth}|p{0.24\textwidth}|}
\hline
\textbf{ Original Comment } & \textbf{ Modification Scope } & \textbf{ Key Revision Excerpt } & \textbf{ Response Explanation } \\
\hline
\endfirsthead
\hline
\textbf{ Original Comment } & \textbf{ Modification Scope } & \textbf{ Key Revision Excerpt } & \textbf{ Response Explanation } \\
\hline
\endhead
\textbf{ Comment: } Please clarify why the transformer model outperforms the CNN baseline in Table 2. In the same response, please also indicate exactly where in the manuscript the supporting evidence is discussed, because at present I can see the gain but not the mechanism-focused explanation behind it. & \texttt{ main.tex::Results::paragraph below Table 2 } & {\color{blue}\ttfamily The transformer improves performance because global-context modeling reduces confusion between visually similar disease patterns while retaining discriminative lesion structure. The supporting evidence for this mechanism-based explanation is provided in the paragraph immediately below Table 2.} & We addressed this combined thread in a single point-to-point row. The revised Results paragraph now explains the mechanism behind the gain and explicitly points readers to the paragraph immediately below Table 2 as the supporting evidence location. \\
\hline
\end{longtable}


\end{document}
